\section{Exact data model}
\label{app:data-model}

This appendix records the finite data conventions used throughout the
native graph, quotient, automorphism, and isomorphism certificates.

\subsection{Vertex ordering}

The native graph uses a fixed ordered vertex list
\[
V(G60)=\{v_0,v_1,\ldots,v_{59}\}.
\]

The symbols $v_i$ denote positions in the canonical exported ordering.
Any alternative labeling gives an isomorphic graph, but all permutation
arrays and certificate digests in this paper refer to this fixed order.

A native permutation
\[
g\in\operatorname{Sym}(V(G60))
\]
is encoded by the length-$60$ tuple
\[
\bigl(g(0),g(1),\ldots,g(59)\bigr).
\]

The identity permutation is therefore
\[
(0,1,\ldots,59).
\]

\subsection{Permutation composition}

Permutation multiplication is function composition. For permutations
$g$ and $h$, the product
\[
gh
\]
means
\[
(gh)(v)=g(h(v)).
\]

In array form,
\[
(gh)[i]=g[h[i]].
\]

All group calculations, generator traversals, multiplication checks, and
mapping certificates use this convention.

\subsection{Native edge representation}

Each undirected native edge is stored as an ordered integer pair
\[
(i,j)
\]
with
\[
i<j.
\]

The complete edge set is sorted lexicographically. This gives a
deterministic serialization of the graph.

A permutation $g$ is accepted as a graph automorphism precisely when
\[
\{\min(g(i),g(j)),\max(g(i),g(j))\}
\]
belongs to the native edge set for every native edge
\[
\{i,j\}\in E(G60).
\]

Because $g$ is a permutation and the edge count is fixed, edge
preservation in one direction is sufficient to establish equality of
the image edge set with the native edge set.

\subsection{Deck involution and fibers}

The canonical deck involution is stored as a length-$60$ permutation
array
\[
a.
\]

It satisfies
\[
a^2=1
\]
and
\[
a(i)\neq i
\]
for every native vertex $i$.

The thirty quotient fibers are the two-element sets
\[
\{i,a(i)\}.
\]

For deterministic storage, each fiber is represented by the ordered pair
\[
(\min(i,a(i)),\max(i,a(i))),
\]
and the complete fiber list is sorted lexicographically.

\subsection{Quotient ordering}

The quotient graph $G30$ uses the sorted fiber list as its canonical
vertex ordering.

Thus quotient vertex $q_k$ corresponds to the $k$th ordered native
fiber.

The projection map
\[
\pi\colon V(G60)\longrightarrow V(G30)
\]
is stored as a length-$60$ array assigning each native vertex its
quotient index.

A quotient permutation is encoded as a length-$30$ tuple in this fixed
quotient ordering.

\subsection{Lift encoding}

For each quotient automorphism
\[
g\in\operatorname{Aut}(G30),
\]
one canonical native lift
\[
\widetilde g
\]
is selected by the deterministic lift construction.

The second lift is
\[
a\widetilde g.
\]

The lift certificate records:

\begin{itemize}
\item the quotient permutation;
\item the canonical native lift;
\item the deck-composed lift;
\item the native edge-preservation result for each lift;
\item the projection compatibility check.
\end{itemize}

Projection compatibility is the equality
\[
\pi\circ\widetilde g=g\circ\pi
\]
tested on all sixty native vertices.

\subsection{Model-group encoding}

Elements of $S_5$ are stored as permutations of a fixed five-point
ordering.

Elements of $D_8$ are stored in normal form
\[
r^i s^j,
\qquad
0\leq i<4,
\qquad
j\in\{0,1\}.
\]

The multiplication law follows from
\[
srs=r^{-1}.
\]

A model element is a pair
\[
(\sigma,r^i s^j)
\]
satisfying
\[
\operatorname{sgn}(\sigma)
=
\chi_{\mathrm{V4}}(r^i s^j).
\]

The complete model group is sorted by the chosen $S_5$ permutation
encoding followed by the dihedral normal-form coordinates.

\subsection{Canonical serialization}

Machine-readable certificates use deterministic key ordering and
deterministic list ordering.

The explicit-isomorphism digest is computed from a canonical
serialization containing:

\begin{itemize}
\item each native permutation tuple;
\item its corresponding $S_5$ coordinate;
\item its corresponding $D_8$ coordinate;
\item the complete mapping order.
\end{itemize}

No timestamps, filesystem paths, or nondeterministic iteration order are
included in the hashed payload.

This makes the digest stable under reproduction on a different machine,
provided the same exact native labeling and mapping are used.
